\section{Gauss 定理和 Stokes 定理}

\begin{problem}{曲面积分}{Surface-Integrals}
    计算下列曲面积分：
    \begin{enumerate}
        \item $\iint_S (x+1) \d y \d z + y \d z \d x + (xy + z) \d x \d y$，其中 $S$ 是以 $O(0,0,0)$、$A(1,0,0)$、$B(0,1,0)$、$C(0,0,1)$ 为顶点的四面体的外表面；
        \item $\iint_S xy \d y \d z + yz \d z \d x + zx \d x \d y$，其中 $S$ 是由 $x=0$、$y=0$、$z=0$、$x+y+z=1$ 所围成的四面体的外侧表面；
        \item $\iint_S x^2 \d y \d z + y^2 \d z \d x + z^2 \d x \d y$，其中 $S$ 是球面 $(x-a)^2 + (y-b)^2 + (z-c)^2 = R^2$ 的外侧；
        \item $\iint_S xy^2 \d y \d z + yz^2 \d z \d x + zx^2 \d x \d y$，其中 $S$ 是球面 $x^2 + y^2 + z^2 = z$ 的外侧；
        \item $\iint_S (x-z) \d y \d z + (y-x) \d z \d x + (z-y) \d x \d y$，其中 $S$ 是旋转抛物面 $z = x^2 + y^2$ ($0 \le z \le 1$) 的下侧；
        \item $\iint_S (y^2+z^2) \d y \d z + (z^2+x^2) \d z \d x + (x^2+y^2) \d x \d y$，其中 $S$ 是上半球面 $x^2 + y^2 + z^2 = a^2$ ($z \ge 0$) 的上侧．
    \end{enumerate}
\end{problem}

\begin{solution}
    \ansitem{1}

    记四面体 $OABC$ 所围成的区域为 $V$．根据高斯公式，有
    \begin{equation}
        \label{eq:part1-1}
        \frac{\partial}{\partial x}(x+1) + \frac{\partial}{\partial y}(y) + \frac{\partial}{\partial z}(xy+z) = 1 + 1 + 1 = 3.
    \end{equation}
    因此，积分值为
    \begin{equation}
        \label{eq:part1-2}
        I = 3 \iiint_V \d x \d y \d z.
    \end{equation}
    由于四面体 $V$ 的体积为
    \begin{equation}
        \label{eq:part1-3}
        \operatorname{Vol}(V) = \frac{1}{6} \times 1 \times 1 \times 1 = \frac{1}{6},
    \end{equation}
    故有
    \begin{equation}
        \label{eq:part1-4}
        I = \boxed{\frac{1}{2}}.
    \end{equation}

    \ansitem{2}

    记四面体所围成的区域为 $V$．根据高斯公式，有
    \begin{equation}
        \label{eq:part2-1}
        I = \iiint_V (x + y + z) \d x \d y \d z.
    \end{equation}
    利用对称性，可得
    \begin{equation}
        \label{eq:part2-2}
        I = 3 \iiint_V x \d x \d y \d z.
    \end{equation}
    其中，三重积分计算如下：
    \begin{equation}
        \label{eq:part2-3}
        \begin{aligned}
            \iiint_V x \d x \d y \d z & = \int_0^1 x \d x \int_0^{1-x} \d y \int_0^{1-x-y} \d z \\
                                      & = \int_0^1 x \d x \int_0^{1-x} (1 - x - y) \d y         \\
                                      & = \frac{1}{2} \int_0^1 x (1 - x)^2 \d x                 \\
                                      & = \frac{1}{2} \int_0^1 (x - 2x^2 + x^3) \d x            \\
                                      & = \frac{1}{24}.
        \end{aligned}
    \end{equation}
    从而有
    \begin{equation}
        \label{eq:part2-4}
        I = \boxed{\frac{1}{8}}.
    \end{equation}

    \ansitem{3}

    记球体为 $V$．根据高斯公式，有
    \begin{equation}
        \label{eq:part3-1}
        I = \iiint_V (2x + 2y + 2z) \d x \d y \d z.
    \end{equation}
    利用形心公式，由于球心位于 $(a, b, c)$，故有
    \begin{equation}
        \label{eq:part3-2}
        \iiint_V x \d x \d y \d z = a \operatorname{Vol}(V), \quad \iiint_V y \d x \d y \d z = b \operatorname{Vol}(V), \quad \iiint_V z \d x \d y \d z = c \operatorname{Vol}(V).
    \end{equation}
    其中，球体体积为
    \begin{equation}
        \label{eq:part3-3}
        \operatorname{Vol}(V) = \frac{4}{3} \uppi R^3.
    \end{equation}
    由此可得
    \begin{equation}
        \label{eq:part3-4}
        I = \boxed{\frac{8}{3} \uppi (a + b + c) R^3}.
    \end{equation}

    \ansitem{4}

    记球体为 $V$．根据高斯公式，有
    \begin{equation}
        \label{eq:part4-1}
        I = \iiint_V (x^2 + y^2 + z^2) \d x \d y \d z.
    \end{equation}
    引入球面坐标系
    \begin{equation}
        \label{eq:part4-2}
        x = r \sin \theta \cos \varphi, \quad y = r \sin \theta \sin \varphi, \quad z = r \cos \theta.
    \end{equation}
    此时，积分区域 $V$ 可表示为 $0 \le \varphi \le 2\uppi$、$0 \le \theta \le \frac{\uppi}{2}$、$0 \le r \le \cos \theta$．从而有
    \begin{equation}
        \label{eq:part4-3}
        \begin{aligned}
            I & = \int_0^{2\uppi} \d \varphi \int_0^{\uppi/2} \d \theta \int_0^{\cos \theta} r^4 \sin \theta \d r \\
              & = 2 \uppi \int_0^{\uppi/2} \frac{\cos^5 \theta}{5} \sin \theta \d \theta                          \\
              & = \frac{2 \uppi}{5} \left[ -\frac{\cos^6 \theta}{6} \right]_0^{\uppi/2}                           \\
              & = \frac{\uppi}{15}.
        \end{aligned}
    \end{equation}
    故有
    \begin{equation}
        \label{eq:part4-4}
        I = \boxed{\frac{\uppi}{15}}.
    \end{equation}

    \ansitem{5}

    曲面 $S$ 的投影区域为 $D = \{(x, y) \mid x^2 + y^2 \le 1\}$，且 $S$ 取下侧，故其法向量的 $z$ 分量为负．将曲面积分转化为投影区域 $D$ 上的二重积分：
    \begin{equation}
        \label{eq:part5-1}
        I = \iint_D \left[ (x - z) \frac{\partial z}{\partial x} + (y - x) \frac{\partial z}{\partial y} - (z - y) \right] \d x \d y.
    \end{equation}
    将 $z = x^2 + y^2$、$\frac{\partial z}{\partial x} = 2x$ 和 $\frac{\partial z}{\partial y} = 2y$ 代入，化简被积函数可得
    \begin{equation}
        \label{eq:part5-2}
        \begin{aligned}
            I & = \iint_D \left[ (x - x^2 - y^2)(2x) + (y - x)(2y) - (x^2 + y^2 - y) \right] \d x \d y \\
              & = \iint_D (x^2 + y^2 + y - 2x^3 - 2xy^2 - 2xy) \d x \d y.
        \end{aligned}
    \end{equation}
    利用积分区域 $D$ 的对称性，奇函数项的积分为零，即
    \begin{equation}
        \label{eq:part5-3}
        \iint_D y \d x \d y = 0, \quad \iint_D x^3 \d x \d y = 0, \quad \iint_D xy^2 \d x \d y = 0, \quad \iint_D xy \d x \d y = 0.
    \end{equation}
    因此，积分化为
    \begin{equation}
        \label{eq:part5-4}
        I = \iint_D (x^2 + y^2) \d x \d y.
    \end{equation}
    引入极坐标系，可得
    \begin{equation}
        \label{eq:part5-5}
        I = \int_0^{2\uppi} \d \theta \int_0^1 r^3 \d r = \frac{\uppi}{2}.
    \end{equation}
    故有
    \begin{equation}
        \label{eq:part5-6}
        I = \boxed{\frac{\uppi}{2}}.
    \end{equation}

    \ansitem{6}

    为了利用高斯公式，添加辅助面 $S_1 = \{(x, y, z) \mid x^2 + y^2 \le a^2, z = 0\}$，其方向取下侧．记 $S_1$ 与 $S$ 围成的上半球体为 $V$．根据高斯公式，有
    \begin{equation}
        \label{eq:part6-1}
        \iint_S + \iint_{S_1} = \iiint_V \left[ \frac{\partial}{\partial x}(y^2+z^2) + \frac{\partial}{\partial y}(z^2+x^2) + \frac{\partial}{\partial z}(x^2+y^2) \right] \d x \d y \d z = 0.
    \end{equation}
    因此，所求积分为
    \begin{equation}
        \label{eq:part6-2}
        I = -\iint_{S_1} (y^2+z^2) \d y \d z + (z^2+x^2) \d z \d x + (x^2+y^2) \d x \d y.
    \end{equation}
    在 $S_1$ 上，有 $z = 0$，且其法向量为 $\symbfit{n} = (0, 0, -1)$，故 $\d y \d z = 0$、$\d z \d x = 0$ 且 $\d x \d y = -\d x \d y$．代入上式得
    \begin{equation}
        \label{eq:part6-3}
        I = \iint_{x^2+y^2 \le a^2} (x^2 + y^2) \d x \d y.
    \end{equation}
    引入极坐标系，可得
    \begin{equation}
        \label{eq:part6-4}
        I = \int_0^{2\uppi} \d \theta \int_0^a r^3 \d r = \frac{\uppi a^4}{2}.
    \end{equation}
    故有
    \begin{equation}
        \label{eq:part6-5}
        I = \boxed{\frac{\uppi a^4}{2}}.
    \end{equation}
\end{solution}

\begin{problem}{引力场的通量}{Gravitational-Flux}
    求引力场 $\symbfit{F} = -km\frac{\symbfit{r}}{r^3}$ 通过下列闭曲面外侧的通量：
    \begin{enumerate}
        \item 空间中任一包围质量 $m$（在原点）的闭曲面；
        \item 空间中任一不包围质量 $m$ 的闭曲面；
        \item 质量 $m$ 在光滑的闭曲面上．
    \end{enumerate}
\end{problem}

\begin{solution}
    \ansitem{1} 通量为 $-4\uppi km$．

    引入直角坐标系，设原点位于质量 $m$ 处，即 $\symbfit{r} = (x, y, z)$，且 $r = \sqrt{x^2+y^2+z^2}$．在 $r \neq 0$ 时，计算引力场 $\symbfit{F}$ 的散度：
    \begin{equation}
        \label{eq:div-F-1}
        \begin{aligned}
            \operatorname{div} \symbfit{F} & = \nabla \cdot \left( -km \frac{\symbfit{r}}{r^3} \right)                                             \\
                                           & = -km \left[ \frac{\nabla \cdot \symbfit{r}}{r^3} + \symbfit{r} \cdot \nabla (r^{-3}) \right]         \\
                                           & = -km \left[ \frac{3}{r^3} + \symbfit{r} \cdot \left( -3 r^{-4} \frac{\symbfit{r}}{r} \right) \right] \\
                                           & = -km \left( \frac{3}{r^3} - \frac{3 r^2}{r^5} \right)                                                \\
                                           & = 0.
        \end{aligned}
    \end{equation}
    设闭曲面 $S$ 包围原点，取以原点为球心、$\varepsilon > 0$ 为半径的足够小球面 $S_\varepsilon = \{ (x, y, z) \mid x^2 + y^2 + z^2 = \varepsilon^2 \}$，使得 $S_\varepsilon$ 被包围在 $S$ 内部．记由 $S$ 与 $S_\varepsilon$ 围成的闭区域为 $V_\varepsilon$．由于在 $V_\varepsilon$ 上 $\operatorname{div} \symbfit{F} = 0$，根据高斯公式可得：
    \begin{equation}
        \label{eq:gauss-1-1}
        \iint_{S} \symbfit{F} \cdot \d \symbfit{S} - \iint_{S_\varepsilon} \symbfit{F} \cdot \d \symbfit{S} = \iiint_{V_\varepsilon} \operatorname{div} \symbfit{F} \d V = 0.
    \end{equation}
    由此可知：
    \begin{equation}
        \label{eq:relation-1}
        \iint_{S} \symbfit{F} \cdot \d \symbfit{S} = \iint_{S_\varepsilon} \symbfit{F} \cdot \d \symbfit{S}.
    \end{equation}
    在 $S_\varepsilon$ 上，外法线单位向量为 $\symbfit{n} = \frac{\symbfit{r}}{\varepsilon}$，此时有：
    \begin{equation}
        \label{eq:integrand-1}
        \symbfit{F} \cdot \symbfit{n} = -km \frac{\symbfit{r}}{\varepsilon^3} \cdot \frac{\symbfit{r}}{\varepsilon} = -\frac{km}{\varepsilon^2}.
    \end{equation}
    因而通量计算为：
    \begin{equation}
        \label{eq:flux-calc-1}
        \iint_{S_\varepsilon} \symbfit{F} \cdot \d \symbfit{S} = \iint_{S_\varepsilon} \left( -\frac{km}{\varepsilon^2} \right) \d S = -\frac{km}{\varepsilon^2} (4\uppi \varepsilon^2) = -4\uppi km.
    \end{equation}
    最终结果为：
    \begin{equation}
        \label{eq:final-ans-1}
        \varphi = \boxed{-4\uppi km}.
    \end{equation}

    \ansitem{2} 通量为 $0$．

    由于闭曲面 $S$ 不包围原点（即不包围质量 $m$），且 $\symbfit{F}$ 在 $S$ 围成的闭区域 $V$ 上连续可微．根据 \zcref{eq:div-F-1}，有 $\operatorname{div} \symbfit{F} = 0$．由高斯公式，可得通量为：
    \begin{equation}
        \label{eq:gauss-2}
        \varphi = \iint_S \symbfit{F} \cdot \d \symbfit{S} = \iiint_V \operatorname{div} \symbfit{F} \d V = 0.
    \end{equation}
    最终结果为：
    \begin{equation}
        \label{eq:final-ans-2}
        \varphi = \boxed{0}.
    \end{equation}

    \ansitem{3} 通量为 $-2\uppi km$．

    由于原点位于光滑的闭曲面 $S$ 上，取以原点为球心、$\varepsilon > 0$ 为半径的小球 $B_\varepsilon$，其边界为 $S_\varepsilon$．记曲面 $S$ 在 $B_\varepsilon$ 外部的部分为 $S \setminus B_\varepsilon$，而 $S_\varepsilon$ 落在 $S$ 内部的球面部分为 $S_{\varepsilon, \text{in}}$．记 $S$ 与 $S_{\varepsilon, \text{in}}$ 共同围成的区域为 $V_\varepsilon$，此时原点在 $V_\varepsilon$ 外部．根据高斯公式有：
    \begin{equation}
        \label{eq:gauss-3}
        \iint_{S \setminus B_\varepsilon} \symbfit{F} \cdot \d \symbfit{S} - \iint_{S_{\varepsilon, \text{in}}} \symbfit{F} \cdot \symbfit{n} \d S = \iiint_{V_\varepsilon} \operatorname{div} \symbfit{F} \d V = 0.
    \end{equation}
    其中 $\symbfit{n}$ 为 $S_\varepsilon$ 的外法线向量．因此有：
    \begin{equation}
        \label{eq:relation-3}
        \iint_{S \setminus B_\varepsilon} \symbfit{F} \cdot \d \symbfit{S} = \iint_{S_{\varepsilon, \text{in}}} \symbfit{F} \cdot \symbfit{n} \d S.
    \end{equation}
    由于 $S$ 在原点处是光滑的，当 $\varepsilon \to 0^+$ 时，$S_{\varepsilon, \text{in}}$ 趋近于一个半球面，其立体角趋近于 $2\uppi$．因此，该部分球面面积满足：
    \begin{equation}
        \label{eq:area-3}
        \operatorname{Area}(S_{\varepsilon, \text{in}}) \sim 2\uppi \varepsilon^2,
        \quad \varepsilon \to 0^+.
    \end{equation}
    代入 $\symbfit{F} \cdot \symbfit{n} = -\frac{km}{\varepsilon^2}$，取极限 $\varepsilon \to 0^+$ 可得：
    \begin{equation}
        \label{eq:limit-3}
        \begin{aligned}
            \varphi & = \lim_{\varepsilon \to 0^+} \iint_{S \setminus B_\varepsilon} \symbfit{F} \cdot \d \symbfit{S}                             \\
                    & = \lim_{\varepsilon \to 0^+} \iint_{S_{\varepsilon, \text{in}}} \left( -\frac{km}{\varepsilon^2} \right) \d S               \\
                    & = \lim_{\varepsilon \to 0^+} \left[ -\frac{km}{\varepsilon^2} \cdot \operatorname{Area}(S_{\varepsilon, \text{in}}) \right] \\
                    & = -2\uppi km.
        \end{aligned}
    \end{equation}
    最终结果为：
    \begin{equation}
        \label{eq:final-ans-3}
        \varphi = \boxed{-2\uppi km}.
    \end{equation}
\end{solution}

\begin{problem}{高斯公式与立体角}{Gauss-Formula-Solid-Angle}
    设区域 $V$ 是由曲面 $x^2+y^2-\frac{z^2}{2}=1$ 及平面 $z=1$、$z=-1$ 围成，$S$ 为 $V$ 的全表面外侧，又设 $\symbfit{V} = (x^2+y^2+z^2)^{-\frac{3}{2}}(x\symbfit{i}+y\symbfit{j}+z\symbfit{k})$．求积分
    \begin{equation}
        \label{eq:integral-def-1}
        \iint_S \frac{x \d y \d z + y \d z \d x + z \d x \d y}{\sqrt{(x^2+y^2+z^2)^3}}.
    \end{equation}
\end{problem}

\begin{solution}
    记所求积分为 $I$．设位置向量为 $\symbfit{r} = x\symbfit{i} + y\symbfit{j} + z\symbfit{k}$，其模长为 $r = \sqrt{x^2+y^2+z^2}$．所求积分实际上是向量场 $\symbfit{V} = \frac{\symbfit{r}}{r^3}$ 穿过闭曲面 $S$ 外侧的通量．

    首先判断原点 $O(0,0,0)$ 与区域 $V$ 的位置关系．将原点坐标代入，满足 $0^2 + 0^2 - \frac{0^2}{2} = 0 < 1$，且介于两平面 $z=-1$ 与 $z=1$ 之间．因此，原点 $O$ 严格位于闭区域 $V$ 的内部．

    对于 $r \neq 0$，计算向量场 $\symbfit{V}$ 的散度：
    \begin{equation}
        \label{eq:div-V}
        \begin{aligned}
            \operatorname{div} \symbfit{V} & = \frac{\partial}{\partial x}\left(\frac{x}{r^3}\right) + \frac{\partial}{\partial y}\left(\frac{y}{r^3}\right) + \frac{\partial}{\partial z}\left(\frac{z}{r^3}\right) \\
                                           & = \left(\frac{1}{r^3} - \frac{3x^2}{r^5}\right) + \left(\frac{1}{r^3} - \frac{3y^2}{r^5}\right) + \left(\frac{1}{r^3} - \frac{3z^2}{r^5}\right)                         \\
                                           & = \frac{3}{r^3} - \frac{3(x^2+y^2+z^2)}{r^5}                                                                                                                            \\
                                           & = 0.
        \end{aligned}
    \end{equation}

    由于原点 $O$ 是该向量场的奇点，我们在区域 $V$ 内部取一个以原点为球心、半径为 $\varepsilon > 0$ 且足够小的球面 $S_\varepsilon \colon x^2+y^2+z^2 = \varepsilon^2$，并取其外侧．记由 $S$ 与 $S_\varepsilon$ 所围成的闭区域为 $V_\varepsilon$．在 $V_\varepsilon$ 上，向量场 $\symbfit{V}$ 连续可微且散度处处为零．

    根据高斯公式，有
    \begin{equation}
        \label{eq:gauss-theorem}
        \iint_S \symbfit{V} \cdot \d\symbfit{S} - \iint_{S_\varepsilon} \symbfit{V} \cdot \d\symbfit{S} = \iiint_{V_\varepsilon} \operatorname{div} \symbfit{V} \d V = 0.
    \end{equation}
    由此可得
    \begin{equation}
        \label{eq:equal-integrals}
        I = \iint_{S} \symbfit{V} \cdot \d\symbfit{S} = \iint_{S_\varepsilon} \symbfit{V} \cdot \d\symbfit{S}.
    \end{equation}

    在球面 $S_\varepsilon$ 上，外法线单位向量为 $\symbfit{n} = \frac{\symbfit{r}}{\varepsilon}$，且在曲面上 $r = \varepsilon$．计算向量场与法线向量的点积：
    \begin{equation}
        \label{eq:integrand-eval}
        \symbfit{V} \cdot \symbfit{n} = \frac{\symbfit{r}}{\varepsilon^3} \cdot \frac{\symbfit{r}}{\varepsilon} = \frac{r^2}{\varepsilon^4} = \frac{\varepsilon^2}{\varepsilon^4} = \frac{1}{\varepsilon^2}.
    \end{equation}
    将 \zcref{eq:integrand-eval} 代入 \zcref{eq:equal-integrals} 中进行积分，即可得到
    \begin{equation}
        \label{eq:final-eval}
        I = \iint_{S_\varepsilon} \frac{1}{\varepsilon^2} \d S = \frac{1}{\varepsilon^2} \operatorname{Area}(S_\varepsilon) = \frac{1}{\varepsilon^2} \left(4\uppi \varepsilon^2\right) = 4\uppi.
    \end{equation}

    最终计算结果为
    \begin{equation}
        \label{eq:result}
        \boxed{4\uppi}.
    \end{equation}
\end{solution}

\begin{problem}{高斯公式与常微分方程}{Gauss-Formula-ODE}
    设对于半空间 $x > 0$ 内任意的光滑有向封闭曲面 $S$，都有
    \begin{equation}
        \label{eq:surface-integral}
        \oiint_S x f(x) \d y \d z - x y f(x) \d z \d x - \e^{2x} z \d x \d y = 0,
    \end{equation}
    其中函数 $f(x)$ 在 $(0, +\infty)$ 内具有连续的一阶导数，且 $\lim_{x \to 0^+} f(x) = 1$，求 $f(x)$．
\end{problem}

\begin{solution}
    记被积向量场为 $\symbfit{F}(x, y, z) = P \symbfit{i} + Q \symbfit{j} + R \symbfit{k}$，其中
    \begin{equation}
        \label{eq:vector-field}
        P = x f(x), \quad Q = -x y f(x), \quad R = -\e^{2x} z.
    \end{equation}
    因为对于半空间 $x > 0$ 内任意的光滑有向封闭曲面 $S$，积分为零，由高斯公式可知，向量场 $\symbfit{F}$ 的散度在 $x > 0$ 内恒为零，即
    \begin{equation}
        \label{eq:div-zero}
        \operatorname{div} \symbfit{F} = \frac{\partial P}{\partial x} + \frac{\partial Q}{\partial y} + \frac{\partial R}{\partial z} = 0.
    \end{equation}
    分别计算各偏导数可得
    \begin{equation}
        \label{eq:partials}
        \begin{aligned}
            \frac{\partial P}{\partial x} & = f(x) + x f'(x), \\
            \frac{\partial Q}{\partial y} & = -x f(x),        \\
            \frac{\partial R}{\partial z} & = -\e^{2x}.
        \end{aligned}
    \end{equation}
    代入 \zcref{eq:div-zero} 中，得到关于 $f(x)$ 的一阶线性常微分方程
    \begin{equation}
        x f'(x) + (1 - x) f(x) - \e^{2x} = 0.
    \end{equation}
    将方程标准化为
    \begin{equation}
        \label{eq:ode-standard}
        f'(x) + \left(\frac{1}{x} - 1\right) f(x) = \frac{\e^{2x}}{x}.
    \end{equation}
    一阶线性微分方程的积分因子为
    \begin{equation}
        \label{eq:integrating-factor}
        \mu(x) = \exp \left( \int \left(\frac{1}{x} - 1\right) \d x \right) = \exp (\ln x - x) = x \e^{-x}.
    \end{equation}
    方程 \zcref{eq:ode-standard} 两边同乘积分因子 $\mu(x)$，得
    \begin{equation}
        \frac{\d}{\d x} \left( x \e^{-x} f(x) \right) = \frac{\e^{2x}}{x} \cdot x \e^{-x} = \e^x.
    \end{equation}
    对方程两边积分，得
    \begin{equation}
        \label{eq:integral-sol}
        x \e^{-x} f(x) = \e^x + C,
    \end{equation}
    其中 $C$ 为任意常数．从而解得通解为
    \begin{equation}
        \label{eq:general-sol}
        f(x) = \frac{\e^{2x} + C \e^x}{x}.
    \end{equation}
    根据题意有极限条件 $\lim_{x \to 0^+} f(x) = 1$，即
    \begin{equation}
        \label{eq:limit-cond}
        \lim_{x \to 0^+} \frac{\e^{2x} + C \e^x}{x} = 1.
    \end{equation}
    由于当 $x \to 0^+$ 时，分母趋于 $0$，要使极限存在且为常数，分子必须趋于 $0$，故要求
    \begin{equation}
        \label{eq:limit-numerator}
        \lim_{x \to 0^+} \left( \e^{2x} + C \e^x \right) = 1 + C = 0,
    \end{equation}
    解得 $C = -1$．将其代回进行洛必达法则验证：
    \begin{equation}
        \label{eq:limit-verify}
        \lim_{x \to 0^+} \frac{\e^{2x} - \e^x}{x} = \lim_{x \to 0^+} \frac{2\e^{2x} - \e^x}{1} = 1,
    \end{equation}
    符合题意．因此，所求函数为
    \begin{equation}
        \label{eq:final-ans}
        \boxed{f(x) = \frac{\e^{2x} - \e^x}{x}}.
    \end{equation}
\end{solution}

\begin{problem}{高斯公式与体积}{Gauss-Formula-Volume}
    证明任意光滑闭曲面 $S$ 围成的立体体积可以表成
    \begin{equation}
        V = \frac{1}{3} \oiint_S x \d y \d z + y \d z \d x + z \d x \d y,
    \end{equation}
    其中积分沿 $S$ 的外侧进行．
\end{problem}

\begin{myproof}
    设由光滑闭曲面 $S$ 所包围的立体区域为 $\Omega$．构造向量场
    \begin{equation}
        \label{eq:vec-field}
        \symbfit{F}(x,y,z) = x \symbfit{i} + y \symbfit{j} + z \symbfit{k}.
    \end{equation}
    计算该向量场的散度，可得
    \begin{equation}
        \label{eq:div-calc}
        \operatorname{div} \symbfit{F} = \frac{\partial x}{\partial x} + \frac{\partial y}{\partial y} + \frac{\partial z}{\partial z} = 1 + 1 + 1 = 3.
    \end{equation}
    由于 $S$ 为任意光滑闭曲面，且积分沿其外侧进行，根据高斯公式，有
    \begin{equation}
        \label{eq:gauss-app-1}
        \oiint_S x \d y \d z + y \d z \d x + z \d x \d y = \iiint_\Omega \operatorname{div} \symbfit{F} \d V.
    \end{equation}
    将 \zcref{eq:div-calc} 代入 \zcref{eq:gauss-app-1} 中，得到
    \begin{equation}
        \label{eq:vol-rel}
        \oiint_S x \d y \d z + y \d z \d x + z \d x \d y = \iiint_\Omega 3 \d V = 3 \iiint_\Omega \d V.
    \end{equation}
    又因为立体 $\Omega$ 的体积 $V$ 定义为
    \begin{equation}
        \label{eq:vol-def}
        V = \iiint_\Omega \d V,
    \end{equation}
    所以 \zcref{eq:vol-rel} 可以写为
    \begin{equation}
        \label{eq:final-proof}
        3V = \oiint_S x \d y \d z + y \d z \d x + z \d x \d y.
    \end{equation}
    等式两边同除以 $3$，即可得到所证等式
    \begin{equation}
        \label{eq:result-1}
        \boxed{V = \frac{1}{3} \oiint_S x \d y \d z + y \d z \d x + z \d x \d y}.
    \end{equation}
\end{myproof}

\begin{problem}{阿基米德原理的证明}{Archimedes-Principle}
    证明 Archimedes （阿基米德）原理：物体 $V$ 全部浸入液体中所受的浮力等于与物体同体积的液体的重量．（提示：设液体的密度为常数 $\rho$，给出物体表面每一小块 $\d S$ 所受到的压力，通过积分计算 $\partial V$ 的压力．）
\end{problem}

\begin{myproof}
    建立空间直角坐标系，设 $z$ 轴竖直向上．设液体内表面的某参考高度为 $z_0$，则在液体内部任意一点 $(x, y, z)$ 处的流体静压强为
    \begin{equation}
        \label{eq:pressure}
        p(z) = p_0 + \rho g (z_0 - z),
    \end{equation}
    其中 $\rho$ 为液体密度，$g$ 为重力加速度，$p_0$ 为参考高度处的压强．

    设物体所占的闭区域为 $V$，其表面为光滑闭曲面 $\partial V$．对于曲面上任意一块面积微元 $\d S$，设其外法线单位向量为 $\symbfit{n} = (\cos \alpha, \cos \beta, \cos \gamma)$，则液体施加在该面积微元上的压力元大小为 $p \d S$，方向沿着内法线方向，即
    \begin{equation}
        \label{eq:force-element}
        \d\symbfit{F} = -p \symbfit{n} \d S.
    \end{equation}

    物体所受的流体总压力 $\symbfit{F}$ 可通过对曲面 $\partial V$ 积分得到．将其分解为三个坐标轴方向的分量：
    \begin{equation}
        \label{eq:force-components}
        \begin{aligned}
            F_x & = \oiint_{\partial V} -p \cos \alpha \d S = -\oiint_{\partial V} p \d y \d z, \\
            F_y & = \oiint_{\partial V} -p \cos \beta \d S = -\oiint_{\partial V} p \d z \d x,  \\
            F_z & = \oiint_{\partial V} -p \cos \gamma \d S = -\oiint_{\partial V} p \d x \d y.
        \end{aligned}
    \end{equation}

    根据高斯公式，将曲面积分转化为体积分，可得
    \begin{equation}
        \label{eq:gauss-app-2}
        \begin{aligned}
            F_x & = -\iiint_V \frac{\partial p}{\partial x} \d V, \\
            F_y & = -\iiint_V \frac{\partial p}{\partial y} \d V, \\
            F_z & = -\iiint_V \frac{\partial p}{\partial z} \d V.
        \end{aligned}
    \end{equation}

    由 \zcref{eq:pressure} 可求得压强的偏导数为
    \begin{equation}
        \label{eq:pressure-derivatives}
        \frac{\partial p}{\partial x} = 0, \quad \frac{\partial p}{\partial y} = 0, \quad \frac{\partial p}{\partial z} = -\rho g.
    \end{equation}

    将 \zcref{eq:pressure-derivatives} 代入 \zcref{eq:gauss-app-2} 中，得到总压力在三个方向的分量为
    \begin{equation}
        \label{eq:final-components}
        F_x = 0, \quad F_y = 0, \quad F_z = -\iiint_V (-\rho g) \d V = \rho g \iiint_V \d V = \rho g V.
    \end{equation}

    因此，物体受到的总压力（即浮力）水平分力为零，方向沿 $z$ 轴正方向（竖直向上），大小等于 $\rho g V$．这就证明了物体全部浸没在液体中所受的浮力刚好等于与其同体积的液体的重量，即
    \begin{equation}
        \label{eq:result-2}
        \boxed{F_z = \rho g V}.
    \end{equation}
\end{myproof}

\begin{problem}{常向量场的曲面积分}{Constant-Vector-Field}
    设 $\symbfit{c}$ 是常向量，$S$ 是任意的光滑闭曲面，证明：
    \begin{equation}
        \iint_S \cos(\widehat{\symbfit{c}, \symbfit{n}}) \d S = 0.
    \end{equation}
    其中 $(\widehat{\symbfit{c}, \symbfit{n}})$ 表示向量 $\symbfit{c}$ 与曲面法向量 $\symbfit{n}$ 的夹角．
\end{problem}

\begin{myproof}
    当 $\symbfit{c} = \symbfit{0}$ 时，被积函数为零，等式显然成立．

    当 $\symbfit{c} \neq \symbfit{0}$ 时，设曲面 $S$ 的外法线单位向量为 $\symbfit{n}$，即 $|\symbfit{n}| = 1$．由向量点积的几何定义，常向量 $\symbfit{c}$ 与法向量 $\symbfit{n}$ 夹角的余弦值可表示为
    \begin{equation}
        \label{eq:cos-def}
        \cos(\widehat{\symbfit{c}, \symbfit{n}}) = \frac{\symbfit{c} \cdot \symbfit{n}}{|\symbfit{c}| |\symbfit{n}|} = \frac{\symbfit{c} \cdot \symbfit{n}}{|\symbfit{c}|}.
    \end{equation}

    对 \zcref{eq:cos-def} 在闭曲面 $S$ 上进行积分，得到
    \begin{equation}
        \label{eq:integral-transform}
        \iint_S \cos(\widehat{\symbfit{c}, \symbfit{n}}) \d S = \iint_S \frac{\symbfit{c} \cdot \symbfit{n}}{|\symbfit{c}|} \d S = \frac{1}{|\symbfit{c}|} \iint_S \symbfit{c} \cdot \symbfit{n} \d S.
    \end{equation}

    设常向量 $\symbfit{c} = (c_1, c_2, c_3)$，其中 $c_1, c_2, c_3$ 均为常数．设闭曲面 $S$ 所围成的三维空间区域为 $V$．根据高斯公式，常向量场 $\symbfit{c}$ 穿过闭曲面 $S$ 的通量可转化为在区域 $V$ 上的体积分：
    \begin{equation}
        \label{eq:gauss-formula}
        \iint_S \symbfit{c} \cdot \symbfit{n} \d S = \iiint_V \operatorname{div} \symbfit{c} \d V.
    \end{equation}

    由于 $\symbfit{c}$ 为常向量场，其各个分量对空间坐标的偏导数均为零，故其散度为
    \begin{equation}
        \label{eq:divergence}
        \operatorname{div} \symbfit{c} = \frac{\partial c_1}{\partial x} + \frac{\partial c_2}{\partial y} + \frac{\partial c_3}{\partial z} = 0.
    \end{equation}

    将 \zcref{eq:divergence} 代入 \zcref{eq:gauss-formula} 中，可得
    \begin{equation}
        \label{eq:flux-zero}
        \iint_S \symbfit{c} \cdot \symbfit{n} \d S = \iiint_V 0 \d V = 0.
    \end{equation}

    进一步将 \zcref{eq:flux-zero} 代回 \zcref{eq:integral-transform}，即可得出所证等式
    \begin{equation}
        \label{eq:final-proof-1}
        \boxed{\iint_S \cos(\widehat{\symbfit{c}, \symbfit{n}}) \d S = 0}.
    \end{equation}
\end{myproof}

\begin{problem}{Gauss 公式与 Green 公式}{Gauss-to-Green}
    设 $L$ 是 $Oxy$ 平面上光滑的简单闭曲线，逆时针方向，立体 $V$ 是柱体，它以 $L$ 为准线，以 $L$ 在 $xy$ 平面内所围平面区域 $D$ 为底，侧面是母线平行于 $z$ 轴的柱面，高为 $1$，试写出向量场 $\symbfit{v} = P\symbfit{i} + Q\symbfit{j}$ 在 $V$ 上的 Gauss 公式，并由此来证明 Green 公式．
\end{problem}

\begin{myproof}
    设向量场为 $\symbfit{v}(x,y,z) = P(x,y)\symbfit{i} + Q(x,y)\symbfit{j} + 0\symbfit{k}$．计算其散度可得
    \begin{equation}
        \label{eq:div-v}
        \operatorname{div} \symbfit{v} = \frac{\partial P}{\partial x} + \frac{\partial Q}{\partial y}.
    \end{equation}
    对柱体区域 $V$ 应用高斯公式，有
    \begin{equation}
        \label{eq:gauss-app}
        \iiint_V \left( \frac{\partial P}{\partial x} + \frac{\partial Q}{\partial y} \right) \d V = \oiint_{\partial V} (P \symbfit{i} + Q \symbfit{j}) \cdot \symbfit{n} \d S.
    \end{equation}
    其中 $\partial V$ 为柱体的全表面外侧，$\symbfit{n}$ 为外法线单位向量．

    对于等式 \zcref{eq:gauss-app} 的左端，由于被积函数仅与 $x,y$ 有关，且柱体高为 $1$（即 $0 \le z \le 1$），化为累次积分可得
    \begin{equation}
        \label{eq:lhs-eval}
        \begin{aligned}
            \iiint_V \left( \frac{\partial P}{\partial x} + \frac{\partial Q}{\partial y} \right) \d V & = \iint_D \left( \frac{\partial P}{\partial x} + \frac{\partial Q}{\partial y} \right) \d x \d y \int_0^1 \d z \\
                                                                                                       & = \iint_D \left( \frac{\partial P}{\partial x} + \frac{\partial Q}{\partial y} \right) \d x \d y.
        \end{aligned}
    \end{equation}

    对于等式 \zcref{eq:gauss-app} 的右端，将全表面分为底面 $S_1$、顶面 $S_2$ 和侧面 $S_3$．
    在底面 $S_1$ ($z=0$) 与顶面 $S_2$ ($z=1$) 上，法向量分别为 $-\symbfit{k}$ 与 $\symbfit{k}$，由于向量场 $\symbfit{v}$ 的 $z$ 分量为零，故有
    \begin{equation}
        \label{eq:top-bottom-flux}
        \iint_{S_1} \symbfit{v} \cdot (-\symbfit{k}) \d S = 0, \quad \iint_{S_2} \symbfit{v} \cdot \symbfit{k} \d S = 0.
    \end{equation}
    对于侧面 $S_3$，其法向量平行于 $xy$ 平面．设曲线 $L$ 的弧长参数为 $s$，由于 $L$ 取逆时针方向，其单位切向量为 $\symbfit{\tau} = \left(\frac{\d x}{\d s}, \frac{\d y}{\d s}, 0\right)$．由此可得侧面外法线单位向量为 $\symbfit{n} = \left(\frac{\d y}{\d s}, -\frac{\d x}{\d s}, 0\right)$．侧面上面积微元为 $\d S = \d s \d z$．因此在 $S_3$ 上的积分为
    \begin{equation}
        \label{eq:side-flux}
        \begin{aligned}
            \iint_{S_3} \symbfit{v} \cdot \symbfit{n} \d S & = \iint_{S_3} \left( P \frac{\d y}{\d s} - Q \frac{\d x}{\d s} \right) \d s \d z      \\
                                                           & = \int_0^1 \d z \oint_L \left( P \frac{\d y}{\d s} - Q \frac{\d x}{\d s} \right) \d s \\
                                                           & = \oint_L -Q \d x + P \d y.
        \end{aligned}
    \end{equation}

    综合 \zcref{eq:lhs-eval}、\zcref{eq:top-bottom-flux} 与 \zcref{eq:side-flux}，高斯公式化为
    \begin{equation}
        \label{eq:green-temp}
        \iint_D \left( \frac{\partial P}{\partial x} + \frac{\partial Q}{\partial y} \right) \d x \d y = \oint_L -Q \d x + P \d y.
    \end{equation}
    上式即为 Green 公式的另一种形式．为了得到标准的 Green 公式，作符号替换：令 $P_1 = -Q$，$Q_1 = P$，代入 \zcref{eq:green-temp} 得
    \begin{equation}
        \label{eq:green-standard}
        \iint_D \left( \frac{\partial Q_1}{\partial x} - \frac{\partial P_1}{\partial y} \right) \d x \d y = \oint_L P_1 \d x + Q_1 \d y.
    \end{equation}
    去掉下标即可得到标准的平面 Green 公式，证毕．
    \begin{equation}
        \label{eq:final-proof-2}
        \boxed{\iint_D \left( \frac{\partial Q}{\partial x} - \frac{\partial P}{\partial y} \right) \d x \d y = \oint_L P \d x + Q \d y}.
    \end{equation}
\end{myproof}

\begin{problem}{曲线积分的计算（Stokes 公式）}{Curve-Integrals-Stokes}
    计算下列曲线积分：
    \begin{enumerate}
        \item $\oint_L y \d x + z \d y + x \d z$，$L$ 是顶点为 $A(1,0,0), B(0,1,0), C(0,0,1)$ 的三角形边界，从原点看去，$L$ 沿顺时针方向；
        \item $\oint_L (y-z) \d x + (z-x) \d y + (x-y) \d z$，$L$ 是圆柱面 $x^2+y^2=a^2$ 和平面 $\frac{x}{a}+\frac{z}{h}=1 \ (a>0, h>0)$ 的交线，从 $x$ 轴的正方向看来，$L$ 沿反时针方向；
        \item $\oint_L (y^2-z^2) \d x + (z^2-x^2) \d y + (x^2-y^2) \d z$，$L$ 是平面 $x+y+z=\frac{3}{2}a$ 与立方体 $0 \leqslant x \leqslant a, 0 \leqslant y \leqslant a, 0 \leqslant z \leqslant a$ 表面的交线，从 $z$ 轴正向看来，$L$ 沿反时针方向；
        \item $\oint_L y^2 \d x + xy \d y + xz \d z$，$L$ 是圆柱面 $x^2+y^2=2y$ 与平面 $y=z$ 的交线，从 $z$ 轴正向看来，$L$ 沿逆时针方向；
        \item $\oint_L (y^2-y) \d x + (z^2-z) \d y + (x^2-x) \d z$，$L$ 是球面 $x^2+y^2+z^2=a^2$ 与平面 $x+y+z=0$ 的交线，$L$ 的方向与 $z$ 轴正向成右手系；
        \item $\oint_L (y^2-z^2) \d x + (2z^2-x^2) \d y + (3x^2-y^2) \d z$，其中 $L$ 是平面 $x+y+z=2$ 与柱面 $|x|+|y|=1$ 的交线，从 $z$ 轴正向看去，$L$ 为逆时针方向．
    \end{enumerate}
\end{problem}

\begin{solution}
    \ansitem{1}

    记 $S$ 为平面 $x+y+z=1$ 在第一象限内的三角形区域．由于从原点看去 $L$ 沿顺时针方向，故 $S$ 的法向量指向远离原点侧，其单位法向量为 $\symbfit{n} = \frac{1}{\sqrt{3}}(1, 1, 1)$．记向量场为 $\symbfit{F}$，计算其旋度为
    \begin{equation}
        \operatorname{curl} \symbfit{F} = \left| \begin{matrix} \symbfit{i} & \symbfit{j} & \symbfit{k} \\ \frac{\partial}{\partial x} & \frac{\partial}{\partial y} & \frac{\partial}{\partial z} \\ y & z & x \end{matrix} \right| = (-1, -1, -1).
    \end{equation}
    利用 Stokes 公式，有
    \begin{equation}
        \label{eq:stokes-1}
        \begin{aligned}
            \oint_L \symbfit{F} \cdot \d \symbfit{r} & = \iint_S \operatorname{curl} \symbfit{F} \cdot \symbfit{n} \d S             \\
                                                     & = \iint_S (-1, -1, -1) \cdot \left( \frac{1}{\sqrt{3}}(1, 1, 1) \right) \d S \\
                                                     & = -\sqrt{3} \iint_S \d S.
        \end{aligned}
    \end{equation}
    三角形面积为 $\iint_S \d S = \frac{\sqrt{3}}{2}$，故
    \begin{equation}
        \label{eq:res-1}
        \boxed{-\frac{3}{2}}.
    \end{equation}

    \ansitem{2}

    记 $S$ 为交线 $L$ 在平面 $\frac{x}{a}+\frac{z}{h}=1$ 上围成的椭圆区域．从 $x$ 轴正向看去 $L$ 沿逆时针方向，故单位法向量 $\symbfit{n}$ 的 $x$ 分量应为正，即 $\symbfit{n} = \frac{(h, 0, a)}{\sqrt{h^2+a^2}}$，因此 $S$ 在 $xy$ 平面的投影 $D \colon x^2+y^2 \leqslant a^2$ 的方向为正．
    计算旋度为
    \begin{equation}
        \operatorname{curl} \symbfit{F} = \left| \begin{matrix} \symbfit{i} & \symbfit{j} & \symbfit{k} \\ \frac{\partial}{\partial x} & \frac{\partial}{\partial y} & \frac{\partial}{\partial z} \\ y-z & z-x & x-y \end{matrix} \right| = (-2, -2, -2).
    \end{equation}
    利用 Stokes 公式化为关于 $xy$ 平面的二重积分（取 $z = h - \frac{h}{a}x$）：
    \begin{equation}
        \label{eq:stokes-2}
        \begin{aligned}
            I & = \iint_D \left[ -(-2)\frac{\partial z}{\partial x} - (-2)\frac{\partial z}{\partial y} + (-2) \right] \d x \d y \\
              & = \iint_D \left[ 2 \left( -\frac{h}{a} \right) + 0 - 2 \right] \d x \d y                                         \\
              & = -2 \left( 1 + \frac{h}{a} \right) \iint_D \d x \d y                                                            \\
              & = -2 \left( \frac{a+h}{a} \right) \cdot \uppi a^2.
        \end{aligned}
    \end{equation}
    故
    \begin{equation}
        \label{eq:res-2}
        \boxed{-2\uppi a(a+h)}.
    \end{equation}

    \ansitem{3}

    记 $S$ 为截口多边形（正六边形区域），其所在平面为 $x+y+z = \frac{3}{2}a$．从 $z$ 轴正向看去 $L$ 为逆时针，取上侧，单位法向量为 $\symbfit{n} = \frac{1}{\sqrt{3}}(1, 1, 1)$．
    计算旋度为
    \begin{equation}
        \label{eq:curl-3}
        \operatorname{curl} \symbfit{F} = (-2y-2z, -2z-2x, -2x-2y) = -2(y+z, x+z, x+y).
    \end{equation}
    由 Stokes 公式，且在平面上满足对称性：
    \begin{equation}
        \label{eq:stokes-3}
        \begin{aligned}
            I & = \iint_S -2(y+z+z+x+x+y) \frac{1}{\sqrt{3}} \d S      \\
              & = -\frac{4}{\sqrt{3}} \iint_S (x+y+z) \d S             \\
              & = -\frac{4}{\sqrt{3}} \cdot \frac{3}{2}a \iint_S \d S.
        \end{aligned}
    \end{equation}
    六边形在 $xy$ 平面的投影 $D$ 的面积为 $\frac{3}{4}a^2$，故 $\iint_S \d S = \sqrt{3} \times \frac{3}{4}a^2$．代入得
    \begin{equation}
        \label{eq:res-3}
        \boxed{-\frac{9}{2}a^3}.
    \end{equation}

    \ansitem{4}

    记 $S$ 为交线在平面 $z=y$ 上的部分，其在 $xy$ 平面投影为 $D \colon x^2+(y-1)^2 \leqslant 1$．从 $z$ 轴正向看逆时针，法向量取上侧．
    计算旋度为
    \begin{equation}
        \label{eq:curl-4}
        \operatorname{curl} \symbfit{F} = (0, -z, -y).
    \end{equation}
    转化为二重积分：
    \begin{equation}
        \label{eq:stokes-4}
        \begin{aligned}
            I & = \iint_D \left[ -(0)\frac{\partial z}{\partial x} - (-z)\frac{\partial z}{\partial y} + (-y) \right] \d x \d y \\
              & = \iint_D (z - y) \d x \d y.
        \end{aligned}
    \end{equation}
    由于在曲面上 $z=y$，故积分为
    \begin{equation}
        \label{eq:res-4}
        \boxed{0}.
    \end{equation}

    \ansitem{5}

    记 $S$ 为交线在平面 $x+y+z=0$ 上所围成的圆盘，半径为 $a$．方向与 $z$ 轴成右手系，故单位法向量 $\symbfit{n} = \frac{1}{\sqrt{3}}(1, 1, 1)$．
    计算旋度为
    \begin{equation}
        \label{eq:curl-5}
        \operatorname{curl} \symbfit{F} = (1-2z, 1-2x, 1-2y).
    \end{equation}
    由 Stokes 公式：
    \begin{equation}
        \label{eq:stokes-5}
        \begin{aligned}
            I & = \iint_S (1-2z+1-2x+1-2y) \frac{1}{\sqrt{3}} \d S \\
              & = \frac{1}{\sqrt{3}} \iint_S [3 - 2(x+y+z)] \d S.
        \end{aligned}
    \end{equation}
    因为 $x+y+z=0$，且 $S$ 面积为 $\uppi a^2$，故
    \begin{equation}
        \label{eq:res-5}
        \boxed{\sqrt{3}\uppi a^2}.
    \end{equation}

    \ansitem{6}

    记 $S$ 为截面位于平面 $z = 2-x-y$ 上的部分，投影区域 $D \colon |x|+|y| \leqslant 1$．从 $z$ 轴看为逆时针，取上侧．
    计算旋度为
    \begin{equation}
        \label{eq:curl-6}
        \operatorname{curl} \symbfit{F} = (-2y-4z, -2z-6x, -2x-2y).
    \end{equation}
    转化为二重积分：
    \begin{equation}
        \label{eq:stokes-6}
        \begin{aligned}
            I & = \iint_D \left[ -(-2y-4z)(-1) - (-2z-6x)(-1) + (-2x-2y) \right] \d x \d y \\
              & = \iint_D (-8x - 4y - 6z) \d x \d y.
        \end{aligned}
    \end{equation}
    将 $z = 2-x-y$ 代入，得被积函数为 $-2x + 2y - 12$．由 $D$ 的对称性，奇函数 $x$ 和 $y$ 的积分为零：
    \begin{equation}
        \label{eq:stokes-6-eval}
        I = \iint_D -12 \d x \d y = -12 \times \operatorname{Area}(D) = -12 \times 2.
    \end{equation}
    故
    \begin{equation}
        \label{eq:res-6}
        \boxed{-24}.
    \end{equation}
\end{solution}

\begin{problem}{Stokes 公式与曲面选取}{Stokes-Surface-Choice}
    在积分 $\oint_L x^2 y^3 \d x + \d y + z \d z$ 中，路径 $L$ 是 $Oxy$ 平面上正向的圆 $x^2+y^2=R^2, z=0$；利用 Stokes 公式化曲线积分为以 $L$ 为边界所围区域 $S$ 上的曲面积分．
    \begin{enumerate}
        \item $S$ 取 $Oxy$ 平面上的圆面 $x^2+y^2 \leqslant R^2$；
        \item $S$ 取半球面 $z = \sqrt{R^2-x^2-y^2}$；
    \end{enumerate}
    结果相同吗？
\end{problem}

\begin{solution}
    \ansitem{1}

    记向量场为 $\symbfit{F}(x,y,z) = x^2 y^3 \symbfit{i} + 1 \symbfit{j} + z \symbfit{k}$．首先计算该向量场的旋度：
    \begin{equation}
        \label{eq:curl-F}
        \operatorname{curl} \symbfit{F} = \left| \begin{matrix} \symbfit{i} & \symbfit{j} & \symbfit{k} \\ \frac{\partial}{\partial x} & \frac{\partial}{\partial y} & \frac{\partial}{\partial z} \\ x^2 y^3 & 1 & z \end{matrix} \right| = (0, 0, -3x^2 y^2).
    \end{equation}
    由于路径 $L$ 在 $Oxy$ 平面上取正向，根据右手法则，所取曲面 $S$ 的法向量应指向上方．对于 $Oxy$ 平面上的圆面 $S_1 \colon x^2+y^2 \leqslant R^2, z=0$，其单位法向量为 $\symbfit{n} = \symbfit{k} = (0,0,1)$．
    利用 Stokes 公式，将曲线积分化为在 $S_1$ 上的曲面积分：
    \begin{equation}
        \label{eq:stokes-S1}
        I = \iint_{S_1} \operatorname{curl} \symbfit{F} \cdot \symbfit{n} \d S = \iint_{S_1} (0, 0, -3x^2 y^2) \cdot (0, 0, 1) \d x \d y = \iint_{x^2+y^2 \leqslant R^2} -3x^2 y^2 \d x \d y.
    \end{equation}
    引入极坐标 $x = r \cos \theta, y = r \sin \theta$，计算该二重积分可得
    \begin{equation}
        \label{eq:eval-S1}
        \begin{aligned}
            I & = -3 \int_0^{2\uppi} \cos^2 \theta \sin^2 \theta \d \theta \int_0^R r^4 \cdot r \d r        \\
              & = -3 \int_0^{2\uppi} \frac{1}{4} \sin^2(2\theta) \d \theta \left[ \frac{r^6}{6} \right]_0^R \\
              & = -3 \left( \frac{\uppi}{4} \right) \left( \frac{R^6}{6} \right)                            \\
              & = -\frac{\uppi R^6}{8}.
        \end{aligned}
    \end{equation}
    最终结果为
    \begin{equation}
        \label{eq:res-1-1}
        \boxed{I = -\frac{\uppi R^6}{8}}.
    \end{equation}

    \ansitem{2}

    取半球面 $S_2 \colon z = \sqrt{R^2-x^2-y^2}$，为使其与边界 $L$ 的正向构成右手系，$S_2$ 需取上侧，故其法向量的方向余弦满足 $\cos \gamma > 0$．根据 Stokes 公式，积分为
    \begin{equation}
        \label{eq:stokes-S2}
        I = \iint_{S_2} \operatorname{curl} \symbfit{F} \cdot \d \symbfit{S} = \iint_{S_2} 0 \d y \d z + 0 \d z \d x - 3x^2 y^2 \d x \d y.
    \end{equation}
    由于曲面 $S_2$ 在 $Oxy$ 平面上的投影区域恰好为圆面 $D \colon x^2+y^2 \leqslant R^2$，且 $S_2$ 取上侧（$\d x \d y$ 符号为正），上述第二型曲面积分可直接化为投影平面上的二重积分：
    \begin{equation}
        \label{eq:eval-S2}
        I = \iint_{D} -3x^2 y^2 \d x \d y.
    \end{equation}
    对比 \zcref{eq:stokes-S1} 可见，该积分式与第 (1) 问完全一致．因此，计算结果同样为
    \begin{equation}
        \label{eq:res-2-1}
        \boxed{I = -\frac{\uppi R^6}{8}}.
    \end{equation}
    这验证了在向量场满足条件的情况下，Stokes 公式的积分结果与所选取的跨越边界的曲面无关，只取决于边界曲线本身．
\end{solution}

\begin{problem}{常向量场的环量}{Constant-Vector-Field-Circulation}
    证明：常向量场 $\symbfit{c}$ 沿任意光滑闭曲线的环量都等于 $0$．
\end{problem}

\begin{myproof}
    设 $L$ 为空间中任意一条光滑闭曲线，其参数方程可表示为位置向量 $\symbfit{r} = \symbfit{r}(t)$，其中参数 $t \in [a, b]$．因为 $L$ 为闭曲线，其起点与终点重合，故有
    \begin{equation}
        \label{eq:closed-curve}
        \symbfit{r}(a) = \symbfit{r}(b).
    \end{equation}

    常向量场 $\symbfit{c}$ 沿闭曲线 $L$ 的环量定义为第二类曲线积分：
    \begin{equation}
        \label{eq:circulation-def}
        \Gamma = \oint_L \symbfit{c} \cdot \d\symbfit{r}.
    \end{equation}

    由于 $\symbfit{c}$ 为常向量，其各个分量不随空间坐标的变化而变化，因此可以根据积分的线性性质将其提取到积分号之外，即
    \begin{equation}
        \label{eq:circulation-extract}
        \Gamma = \symbfit{c} \cdot \left( \oint_L \d\symbfit{r} \right).
    \end{equation}

    对于闭曲线 $L$ 上微元 $\d\symbfit{r}$ 的积分，根据微积分基本定理以及 \zcref{eq:closed-curve} 可得
    \begin{equation}
        \label{eq:integral-dr}
        \oint_L \d\symbfit{r} = \int_a^b \frac{\d\symbfit{r}}{\d t} \d t = \symbfit{r}(b) - \symbfit{r}(a) = \symbfit{0}.
    \end{equation}

    将 \zcref{eq:integral-dr} 代入 \zcref{eq:circulation-extract} 中进行计算，可得出该环量的值为
    \begin{equation}
        \label{eq:circulation-calc}
        \Gamma = \symbfit{c} \cdot \symbfit{0} = 0.
    \end{equation}

    综上所述，常向量场沿任意光滑闭曲线的环量恒为 $0$，即
    \begin{equation}
        \label{eq:final-proof-3}
        \boxed{\oint_L \symbfit{c} \cdot \d\symbfit{r} = 0}.
    \end{equation}
\end{myproof}

\begin{problem}{Stokes 公式与环量}{Circulation-Stokes}
    求向量场 $\symbfit{v} = (y^2+z^2)\symbfit{i} + (z^2+x^2)\symbfit{j} + (x^2+y^2)\symbfit{k}$ 沿曲线 $L$ 的环量．$L$ 为 $x^2+y^2+z^2=R^2 \ (z \geqslant 0)$ 与 $x^2+y^2=Rx$ 的交线，从 $x$ 轴正向看来，$L$ 沿反时针方向．
\end{problem}

\begin{solution}
    记向量场为 $\symbfit{v} = P\symbfit{i} + Q\symbfit{j} + R\symbfit{k}$，其中分量分别为
    \begin{equation}
        \label{eq:vector-components}
        P = y^2+z^2, \quad Q = z^2+x^2, \quad R = x^2+y^2.
    \end{equation}
    首先计算向量场 $\symbfit{v}$ 的旋度：
    \begin{equation}
        \label{eq:curl-v}
        \begin{aligned}
            \operatorname{curl} \symbfit{v} & = \left( \frac{\partial R}{\partial y} - \frac{\partial Q}{\partial z} \right)\symbfit{i} + \left( \frac{\partial P}{\partial z} - \frac{\partial R}{\partial x} \right)\symbfit{j} + \left( \frac{\partial Q}{\partial x} - \frac{\partial P}{\partial y} \right)\symbfit{k} \\
                                            & = (2y - 2z)\symbfit{i} + (2z - 2x)\symbfit{j} + (2x - 2y)\symbfit{k}.
        \end{aligned}
    \end{equation}

    取积分曲面 $S$ 为上半球面 $x^2+y^2+z^2=R^2$ 在柱面 $x^2+y^2=Rx$ 内部的部分．根据题意，从 $x$ 轴正向看来（即视线迎着 $x$ 轴正方向），曲线 $L$ 的走向为逆时针方向．由右手定则可知，曲面 $S$ 上的单位法向量应具有正的 $x$ 分量．对于原点在球心的球面外侧，其单位法向量恰为
    \begin{equation}
        \label{eq:normal-vector}
        \symbfit{n} = \frac{1}{R} (x\symbfit{i} + y\symbfit{j} + z\symbfit{k}).
    \end{equation}
    由于在曲面上 $x \geqslant 0$（由 $x^2+y^2=Rx$ 决定），取球面外侧的方向与曲线 $L$ 满足右手系定向要求．

    根据 Stokes 公式，向量场沿曲线 $L$ 的环量可以转化为在曲面 $S$ 上的第一类曲面积分：
    \begin{equation}
        \label{eq:stokes-formula}
        \oint_L \symbfit{v} \cdot \d\symbfit{r} = \iint_S \operatorname{curl} \symbfit{v} \cdot \symbfit{n} \d S.
    \end{equation}
    计算被积函数中的点积：
    \begin{equation}
        \label{eq:dot-product}
        \begin{aligned}
            \operatorname{curl} \symbfit{v} \cdot \symbfit{n} & = \frac{1}{R} \big[ (2y-2z)x + (2z-2x)y + (2x-2y)z \big] \\
                                                              & = \frac{1}{R} (2xy - 2zx + 2zy - 2xy + 2xz - 2yz)        \\
                                                              & = 0.
        \end{aligned}
    \end{equation}

    由于在曲面 $S$ 上处处有 $\operatorname{curl} \symbfit{v} \cdot \symbfit{n} = 0$，将其代入曲面积分中可得
    \begin{equation}
        \label{eq:result-integral}
        \oint_L \symbfit{v} \cdot \d\symbfit{r} = \iint_S 0 \d S = 0.
    \end{equation}

    因此，所求环量为
    \begin{equation}
        \label{eq:final-answer}
        \boxed{0}.
    \end{equation}
\end{solution}

\endinput
